% --- Archon physics formalization source begin ---
% archon:physics
% archon:covers PhyXMiniProblems/problem_phyx_mini_0994.lean
% archon:source-report reports/phyx_mini/problem_phyx_mini_0994.source.json

\chapter{Physics problem phyx\_mini\_0994}
\label{ch:PhyXMiniProblems_problem_phyx_mini_0994}

\paragraph{Problem source.}
\#\# Physical scenario

Point charges \textbackslash{}( q\_1 = +12\textbackslash{}text\{nC\} \textbackslash{}) and \textbackslash{}( q\_2 = -12\textbackslash{}text\{nC\} \textbackslash{}) are \textbackslash{}( 0.100\textbackslash{}text\{m\} \textbackslash{}) apart. (Such pairs of point charges with equal magnitude and opposite sign are called \textbackslash{}textit\{electric dipoles\}.)

\#\# Question

Compute the total field at point \textbackslash{}( c \textbackslash{}).

\#\# Answer choices

- A: A: (6.39 \textbackslash{}times 10\textasciicircum{}3\textbackslash{}

- B: B: (4.9 \textbackslash{}times 10\textasciicircum{}3\textbackslash{}

- C: C: (8.78 \textbackslash{}times 10\textasciicircum{}3\textbackslash{}

- D: D: (5.3 \textbackslash{}times 10\textasciicircum{}3\textbackslash{}

\#\# Figure caption (auxiliary; use the image as primary evidence)

The image is a physics diagram illustrating the electric fields produced by two point charges, \textbackslash{}( q\_1 \textbackslash{}) and \textbackslash{}( q\_2 \textbackslash{}), positioned on a horizontal axis labeled \textbackslash{}( x \textbackslash{}), with a vertical axis labeled \textbackslash{}( y \textbackslash{}).

- **Charge \textbackslash{}( q\_1 \textbackslash{}):** This is a positive charge located 4.0 cm from a point \textbackslash{}( b \textbackslash{}) on the left along the \textbackslash{}( x \textbackslash{})-axis.

- **Charge \textbackslash{}( q\_2 \textbackslash{}):** This is a negative charge located 4.0 cm from a point \textbackslash{}( a \textbackslash{}) on the right, along the \textbackslash{}( x \textbackslash{})-axis.

- **Direction and Magnitude:** 
  - \textbackslash{}( \textbackslash{}vec\{E\}\_b \textbackslash{}) is directed leftward from point \textbackslash{}( b \textbackslash{}).
  - \textbackslash{}( \textbackslash{}vec\{E\}\_a \textbackslash{}) is directed rightward from point \textbackslash{}( a \textbackslash{}).

- **Central point \textbackslash{}( c \textbackslash{}):** Positioned vertically above the midpoint between charges \textbackslash{}( q\_1 \textbackslash{}) and \textbackslash{}( q\_2 \textbackslash{}) at a distance of 13.0 cm along dashed lines forming a triangle. This point has vectors \textbackslash{}( \textbackslash{}vec\{E\}\_1 \textbackslash{}) and \textbackslash{}( \textbackslash{}vec\{E\}\_2 \textbackslash{}) originating from it.
  
- **Angles:** The angles \textbackslash{}( \textbackslash{}alpha \textbackslash{}) are measured from the horizontal to the vectors directed from point \textbackslash{}( c \textbackslash{}) toward both charges.

- **Vectors:**
  - \textbackslash{}( \textbackslash{}vec\{E\}\_1 \textbackslash{}) and \textbackslash{}( \textbackslash{}vec\{E\}\_2 \textbackslash{}) are electric field vectors originating from the central point \textbackslash{}( c \textbackslash{}) and directed downward toward charges \textbackslash{}( q\_1 \textbackslash{}) and \textbackslash{}( q\_2 \textbackslash{}), respectively. 
  - They are at angles \textbackslash{}( \textbackslash{}alpha \textbackslash{}) with respect to the vertical downward direction, combining to give a resultant field \textbackslash{}( \textbackslash{}vec\{E\}\_c \textbackslash{}) directed downward.

- **Distances:** 
  - Points \textbackslash{}( b \textbackslash{}) and \textbackslash{}( a \textbackslash{}) are each 4.0 cm from their respective charges, and points \textbackslash{}( c \textbackslash{}) are 13.0 cm from the charges measured along the dashed lines which form a triangle.

Overall, the image depicts the configuration of electric fields due to two charges and their geometric and vectorial interactions, highlighting the principles of electrostatics.

\#\# Dataset metadata

Electrostatics; Physical Model Grounding Reasoning, Spatial Relation Reasoning

\paragraph{Recorded answer/context.}
B: B: (4.9 \textbackslash{}times 10\textasciicircum{}3\textbackslash{}

\paragraph{Figure/image path.}
/root/proposal\_for\_physic/science-mango/phyx\_mini\_run/phyx\_data/test\_image/994.png

\paragraph{Formalization target.}
create a compiling Lean file with sorry bodies at `PhyXMiniProblems/problem\_phyx\_mini\_0994.lean`. The Lean declarations must preserve the physical quantities, dimensions or dimensional roles, figure labels, governing-law hypotheses, and final relation expressed by this problem.
Use Mathlib/Physlib names found through LeanExplore where available. If a physics API is missing, introduce faithful local abstractions rather than scalar placeholder aliases.

\begin{theorem}[Physics formalization target]
\label{thm:physics:phyx_mini_0994:target}
\lean{PhyXMiniProblems.ProblemPhyXMini0994.problem_phyx_mini_0994}
\uses{def:physics:phyx-mini-0994:phyxminiproblems-problemphyxmini0994-fieldmagnitudeinnewtonspercoulomb, def:physics:phyx-mini-0994:phyxminiproblems-problemphyxmini0994-electricdipolesetup, def:physics:phyx-mini-0994:phyxminiproblems-problemphyxmini0994-resultantfieldatcinnewtonspercoulomb, def:physics:phyx-mini-0994:phyxminiproblems-problemphyxmini0994-matchesstatedelectricdipolescenario, def:physics:phyx-mini-0994:phyxminiproblems-problemphyxmini0994-matchessuppliedelectricdipolefigure, def:physics:phyx-mini-0994:phyxminiproblems-problemphyxmini0994-hasphysicalelectricdipoleparameters, def:physics:phyx-mini-0994:phyxminiproblems-problemphyxmini0994-usesschoolcoulombconstant, def:physics:phyx-mini-0994:phyxminiproblems-problemphyxmini0994-satisfiespointchargecoulombfieldlaw, def:physics:phyx-mini-0994:phyxminiproblems-problemphyxmini0994-satisfieselectricfieldsuperposition, def:physics:phyx-mini-0994:phyxminiproblems-problemphyxmini0994-resultantmagnituderepresentsfieldatc, def:physics:phyx-mini-0994:phyxminiproblems-problemphyxmini0994-recordeddatasetanswer, def:physics:phyx-mini-0994:phyxminiproblems-problemphyxmini0994-isrightward, def:physics:phyx-mini-0994:phyxminiproblems-problemphyxmini0994-roundstonearesthundred, def:physics:phyx-mini-0994:phyxminiproblems-problemphyxmini0994-isuniquenearestdisplayedfieldmagnitude}
The assigned autoformalize agent should translate this physics problem into Lean declarations in the covered file, with theorem and lemma proof bodies written as `by sorry`.
\end{theorem}
\begin{proof}
This is an autoformalization task, not a proof task. Produce statements that can later be proved by the physics prover without weakening the physical model.
\end{proof}

% --- Archon declaration topology begin ---
\section{Declaration topology}
\begin{definition}[Length Quantity]
  \label{def:physics:phyx-mini-0994:phyxminiproblems-problemphyxmini0994-lengthquantity}
  \lean{PhyXMiniProblems.ProblemPhyXMini0994.LengthQuantity}
  A nonnegative, unit-independent physical length
\end{definition}
\begin{proof}
  This is the declared model definition of Length Quantity.
\end{proof}
\begin{definition}[Signed Charge Quantity]
  \label{def:physics:phyx-mini-0994:phyxminiproblems-problemphyxmini0994-signedchargequantity}
  \lean{PhyXMiniProblems.ProblemPhyXMini0994.SignedChargeQuantity}
  A signed, unit-independent physical electric charge
\end{definition}
\begin{proof}
  This is the declared model definition of Signed Charge Quantity.
\end{proof}
\begin{definition}[electric Field Strength Dimension]
  \label{def:physics:phyx-mini-0994:phyxminiproblems-problemphyxmini0994-electricfieldstrengthdimension}
  \lean{PhyXMiniProblems.ProblemPhyXMini0994.electricFieldStrengthDimension}
  The physical dimension M L T⁻² C⁻¹ of electric-field strength
\end{definition}
\begin{proof}
  This is the declared model definition of electric Field Strength Dimension.
\end{proof}
\begin{definition}[Electric Field Magnitude Quantity]
  \label{def:physics:phyx-mini-0994:phyxminiproblems-problemphyxmini0994-electricfieldmagnitudequantity}
  \lean{PhyXMiniProblems.ProblemPhyXMini0994.ElectricFieldMagnitudeQuantity}
  \uses{def:physics:phyx-mini-0994:phyxminiproblems-problemphyxmini0994-electricfieldstrengthdimension}
  A nonnegative physical electric-field magnitude
\end{definition}
\begin{proof}
  This is the declared model definition of Electric Field Magnitude Quantity.
\end{proof}
\begin{definition}[length In Meters]
  \label{def:physics:phyx-mini-0994:phyxminiproblems-problemphyxmini0994-lengthinmeters}
  \lean{PhyXMiniProblems.ProblemPhyXMini0994.lengthInMeters}
  \uses{def:physics:phyx-mini-0994:phyxminiproblems-problemphyxmini0994-lengthquantity}
  Coherent-SI metre readout of a physical length
\end{definition}
\begin{proof}
  This is the declared model definition of length In Meters.
\end{proof}
\begin{definition}[length In Centimeters]
  \label{def:physics:phyx-mini-0994:phyxminiproblems-problemphyxmini0994-lengthincentimeters}
  \lean{PhyXMiniProblems.ProblemPhyXMini0994.lengthInCentimeters}
  \uses{def:physics:phyx-mini-0994:phyxminiproblems-problemphyxmini0994-lengthquantity, def:physics:phyx-mini-0994:phyxminiproblems-problemphyxmini0994-lengthinmeters}
  Centimetre readout used by the dimension labels in the raster
\end{definition}
\begin{proof}
  This is the declared model definition of length In Centimeters.
\end{proof}
\begin{definition}[charge In Coulombs]
  \label{def:physics:phyx-mini-0994:phyxminiproblems-problemphyxmini0994-chargeincoulombs}
  \lean{PhyXMiniProblems.ProblemPhyXMini0994.chargeInCoulombs}
  \uses{def:physics:phyx-mini-0994:phyxminiproblems-problemphyxmini0994-signedchargequantity}
  Coherent-SI coulomb readout of a signed physical charge
\end{definition}
\begin{proof}
  This is the declared model definition of charge In Coulombs.
\end{proof}
\begin{definition}[charge In Nanocoulombs]
  \label{def:physics:phyx-mini-0994:phyxminiproblems-problemphyxmini0994-chargeinnanocoulombs}
  \lean{PhyXMiniProblems.ProblemPhyXMini0994.chargeInNanocoulombs}
  \uses{def:physics:phyx-mini-0994:phyxminiproblems-problemphyxmini0994-signedchargequantity, def:physics:phyx-mini-0994:phyxminiproblems-problemphyxmini0994-chargeincoulombs}
  Nanocoulomb readout used by the two charge labels
\end{definition}
\begin{proof}
  This is the declared model definition of charge In Nanocoulombs.
\end{proof}
\begin{definition}[field Magnitude In Newtons Per Coulomb]
  \label{def:physics:phyx-mini-0994:phyxminiproblems-problemphyxmini0994-fieldmagnitudeinnewtonspercoulomb}
  \lean{PhyXMiniProblems.ProblemPhyXMini0994.fieldMagnitudeInNewtonsPerCoulomb}
  \uses{def:physics:phyx-mini-0994:phyxminiproblems-problemphyxmini0994-electricfieldmagnitudequantity}
  Coherent-SI readout of a field magnitude in newtons per coulomb
\end{definition}
\begin{proof}
  This is the declared model definition of field Magnitude In Newtons Per Coulomb.
\end{proof}
\begin{definition}[Source Charge]
  \label{def:physics:phyx-mini-0994:phyxminiproblems-problemphyxmini0994-sourcecharge}
  \lean{PhyXMiniProblems.ProblemPhyXMini0994.SourceCharge}
  The two point charges named in the problem and raster
\end{definition}
\begin{proof}
  This is the declared model definition of Source Charge.
\end{proof}
\begin{definition}[Figure Point]
  \label{def:physics:phyx-mini-0994:phyxminiproblems-problemphyxmini0994-figurepoint}
  \lean{PhyXMiniProblems.ProblemPhyXMini0994.FigurePoint}
  All five labelled points on the horizontal/vertical diagram
\end{definition}
\begin{proof}
  This is the declared model definition of Figure Point.
\end{proof}
\begin{definition}[Figure Charge Sign]
  \label{def:physics:phyx-mini-0994:phyxminiproblems-problemphyxmini0994-figurechargesign}
  \lean{PhyXMiniProblems.ProblemPhyXMini0994.FigureChargeSign}
  The sign glyph printed inside each charge circle
\end{definition}
\begin{proof}
  This is the declared model definition of Figure Charge Sign.
\end{proof}
\begin{definition}[Figure Charge Color]
  \label{def:physics:phyx-mini-0994:phyxminiproblems-problemphyxmini0994-figurechargecolor}
  \lean{PhyXMiniProblems.ProblemPhyXMini0994.FigureChargeColor}
  Fill colours used to distinguish the two charge signs
\end{definition}
\begin{proof}
  This is the declared model definition of Figure Charge Color.
\end{proof}
\begin{definition}[Figure Segment]
  \label{def:physics:phyx-mini-0994:phyxminiproblems-problemphyxmini0994-figuresegment}
  \lean{PhyXMiniProblems.ProblemPhyXMini0994.FigureSegment}
  The five distance segments carrying labels or dashed geometry
\end{definition}
\begin{proof}
  This is the declared model definition of Figure Segment.
\end{proof}
\begin{definition}[Figure Field Arrow]
  \label{def:physics:phyx-mini-0994:phyxminiproblems-problemphyxmini0994-figurefieldarrow}
  \lean{PhyXMiniProblems.ProblemPhyXMini0994.FigureFieldArrow}
  The electric-field arrows explicitly labelled in the supplied image
\end{definition}
\begin{proof}
  This is the declared model definition of Figure Field Arrow.
\end{proof}
\begin{definition}[Figure Arrow Direction]
  \label{def:physics:phyx-mini-0994:phyxminiproblems-problemphyxmini0994-figurearrowdirection}
  \lean{PhyXMiniProblems.ProblemPhyXMini0994.FigureArrowDirection}
  Qualitative directions sufficient to transcribe the raster faithfully
\end{definition}
\begin{proof}
  This is the declared model definition of Figure Arrow Direction.
\end{proof}
\begin{definition}[Alpha Mark]
  \label{def:physics:phyx-mini-0994:phyxminiproblems-problemphyxmini0994-alphamark}
  \lean{PhyXMiniProblems.ProblemPhyXMini0994.AlphaMark}
  The three occurrences of the angle label α in the image
\end{definition}
\begin{proof}
  This is the declared model definition of Alpha Mark.
\end{proof}
\begin{definition}[expected Source Point]
  \label{def:physics:phyx-mini-0994:phyxminiproblems-problemphyxmini0994-expectedsourcepoint}
  \lean{PhyXMiniProblems.ProblemPhyXMini0994.expectedSourcePoint}
  \uses{def:physics:phyx-mini-0994:phyxminiproblems-problemphyxmini0994-sourcecharge, def:physics:phyx-mini-0994:phyxminiproblems-problemphyxmini0994-figurepoint}
  The image point occupied by each physical source charge
\end{definition}
\begin{proof}
  This is the declared model definition of expected Source Point.
\end{proof}
\begin{definition}[expected Charge In Nanocoulombs]
  \label{def:physics:phyx-mini-0994:phyxminiproblems-problemphyxmini0994-expectedchargeinnanocoulombs}
  \lean{PhyXMiniProblems.ProblemPhyXMini0994.expectedChargeInNanocoulombs}
  \uses{def:physics:phyx-mini-0994:phyxminiproblems-problemphyxmini0994-sourcecharge}
  Signed nanocoulomb value printed beside each source
\end{definition}
\begin{proof}
  This is the declared model definition of expected Charge In Nanocoulombs.
\end{proof}
\begin{definition}[expected Charge Sign]
  \label{def:physics:phyx-mini-0994:phyxminiproblems-problemphyxmini0994-expectedchargesign}
  \lean{PhyXMiniProblems.ProblemPhyXMini0994.expectedChargeSign}
  \uses{def:physics:phyx-mini-0994:phyxminiproblems-problemphyxmini0994-sourcecharge, def:physics:phyx-mini-0994:phyxminiproblems-problemphyxmini0994-figurechargesign}
  Sign glyph printed inside each source circle
\end{definition}
\begin{proof}
  This is the declared model definition of expected Charge Sign.
\end{proof}
\begin{definition}[expected Charge Color]
  \label{def:physics:phyx-mini-0994:phyxminiproblems-problemphyxmini0994-expectedchargecolor}
  \lean{PhyXMiniProblems.ProblemPhyXMini0994.expectedChargeColor}
  \uses{def:physics:phyx-mini-0994:phyxminiproblems-problemphyxmini0994-sourcecharge, def:physics:phyx-mini-0994:phyxminiproblems-problemphyxmini0994-figurechargecolor}
  Colour of each source circle in the raster
\end{definition}
\begin{proof}
  This is the declared model definition of expected Charge Color.
\end{proof}
\begin{definition}[segment Start]
  \label{def:physics:phyx-mini-0994:phyxminiproblems-problemphyxmini0994-segmentstart}
  \lean{PhyXMiniProblems.ProblemPhyXMini0994.segmentStart}
  \uses{def:physics:phyx-mini-0994:phyxminiproblems-problemphyxmini0994-figurepoint, def:physics:phyx-mini-0994:phyxminiproblems-problemphyxmini0994-figuresegment}
  First endpoint of each dimensioned or dashed segment
\end{definition}
\begin{proof}
  This is the declared model definition of segment Start.
\end{proof}
\begin{definition}[segment End]
  \label{def:physics:phyx-mini-0994:phyxminiproblems-problemphyxmini0994-segmentend}
  \lean{PhyXMiniProblems.ProblemPhyXMini0994.segmentEnd}
  \uses{def:physics:phyx-mini-0994:phyxminiproblems-problemphyxmini0994-figurepoint, def:physics:phyx-mini-0994:phyxminiproblems-problemphyxmini0994-figuresegment}
  Second endpoint of each dimensioned or dashed segment
\end{definition}
\begin{proof}
  This is the declared model definition of segment End.
\end{proof}
\begin{definition}[expected Segment Length In Centimeters]
  \label{def:physics:phyx-mini-0994:phyxminiproblems-problemphyxmini0994-expectedsegmentlengthincentimeters}
  \lean{PhyXMiniProblems.ProblemPhyXMini0994.expectedSegmentLengthInCentimeters}
  \uses{def:physics:phyx-mini-0994:phyxminiproblems-problemphyxmini0994-figuresegment}
  Centimetre values literally printed on the five figure segments
\end{definition}
\begin{proof}
  This is the declared model definition of expected Segment Length In Centimeters.
\end{proof}
\begin{definition}[expected Arrow Origin]
  \label{def:physics:phyx-mini-0994:phyxminiproblems-problemphyxmini0994-expectedarroworigin}
  \lean{PhyXMiniProblems.ProblemPhyXMini0994.expectedArrowOrigin}
  \uses{def:physics:phyx-mini-0994:phyxminiproblems-problemphyxmini0994-figurepoint, def:physics:phyx-mini-0994:phyxminiproblems-problemphyxmini0994-figurefieldarrow}
  The point from which each labelled arrow is drawn
\end{definition}
\begin{proof}
  This is the declared model definition of expected Arrow Origin.
\end{proof}
\begin{definition}[expected Arrow Direction]
  \label{def:physics:phyx-mini-0994:phyxminiproblems-problemphyxmini0994-expectedarrowdirection}
  \lean{PhyXMiniProblems.ProblemPhyXMini0994.expectedArrowDirection}
  \uses{def:physics:phyx-mini-0994:phyxminiproblems-problemphyxmini0994-figurefieldarrow, def:physics:phyx-mini-0994:phyxminiproblems-problemphyxmini0994-figurearrowdirection}
  Qualitative arrow directions read directly from the primary raster
\end{definition}
\begin{proof}
  This is the declared model definition of expected Arrow Direction.
\end{proof}
\begin{definition}[Electric Dipole Figure]
  \label{def:physics:phyx-mini-0994:phyxminiproblems-problemphyxmini0994-electricdipolefigure}
  \lean{PhyXMiniProblems.ProblemPhyXMini0994.ElectricDipoleFigure}
  \uses{def:physics:phyx-mini-0994:phyxminiproblems-problemphyxmini0994-sourcecharge, def:physics:phyx-mini-0994:phyxminiproblems-problemphyxmini0994-figurepoint, def:physics:phyx-mini-0994:phyxminiproblems-problemphyxmini0994-figurechargesign, def:physics:phyx-mini-0994:phyxminiproblems-problemphyxmini0994-figurechargecolor, def:physics:phyx-mini-0994:phyxminiproblems-problemphyxmini0994-figuresegment, def:physics:phyx-mini-0994:phyxminiproblems-problemphyxmini0994-figurefieldarrow, def:physics:phyx-mini-0994:phyxminiproblems-problemphyxmini0994-figurearrowdirection, def:physics:phyx-mini-0994:phyxminiproblems-problemphyxmini0994-alphamark}
  Literal presentation data transcribed from image 994.png
\end{definition}
\begin{proof}
  This is the declared model definition of Electric Dipole Figure.
\end{proof}
\begin{definition}[position Vector In Meters]
  \label{def:physics:phyx-mini-0994:phyxminiproblems-problemphyxmini0994-positionvectorinmeters}
  \lean{PhyXMiniProblems.ProblemPhyXMini0994.positionVectorInMeters}
  Turn a planar Space 2 point into its explicitly metre-valued vector
\end{definition}
\begin{proof}
  This is the declared model definition of position Vector In Meters.
\end{proof}
\begin{definition}[Electric Dipole Setup]
  \label{def:physics:phyx-mini-0994:phyxminiproblems-problemphyxmini0994-electricdipolesetup}
  \lean{PhyXMiniProblems.ProblemPhyXMini0994.ElectricDipoleSetup}
  \uses{def:physics:phyx-mini-0994:phyxminiproblems-problemphyxmini0994-lengthquantity, def:physics:phyx-mini-0994:phyxminiproblems-problemphyxmini0994-signedchargequantity, def:physics:phyx-mini-0994:phyxminiproblems-problemphyxmini0994-electricfieldmagnitudequantity, def:physics:phyx-mini-0994:phyxminiproblems-problemphyxmini0994-sourcecharge, def:physics:phyx-mini-0994:phyxminiproblems-problemphyxmini0994-figurepoint, def:physics:phyx-mini-0994:phyxminiproblems-problemphyxmini0994-electricdipolefigure}
  Independent physical objects in the dipole setup. In particular, the source fields, resultant field, and resultant magnitude are observables, not local definitions of the requested answer
\end{definition}
\begin{proof}
  This is the declared model definition of Electric Dipole Setup.
\end{proof}
\begin{definition}[source Position]
  \label{def:physics:phyx-mini-0994:phyxminiproblems-problemphyxmini0994-sourceposition}
  \lean{PhyXMiniProblems.ProblemPhyXMini0994.sourcePosition}
  \uses{def:physics:phyx-mini-0994:phyxminiproblems-problemphyxmini0994-sourcecharge, def:physics:phyx-mini-0994:phyxminiproblems-problemphyxmini0994-electricdipolesetup}
  Spatial position occupied by a named source charge
\end{definition}
\begin{proof}
  This is the declared model definition of source Position.
\end{proof}
\begin{definition}[displacement From Source In Meters]
  \label{def:physics:phyx-mini-0994:phyxminiproblems-problemphyxmini0994-displacementfromsourceinmeters}
  \lean{PhyXMiniProblems.ProblemPhyXMini0994.displacementFromSourceInMeters}
  \uses{def:physics:phyx-mini-0994:phyxminiproblems-problemphyxmini0994-sourcecharge, def:physics:phyx-mini-0994:phyxminiproblems-problemphyxmini0994-positionvectorinmeters, def:physics:phyx-mini-0994:phyxminiproblems-problemphyxmini0994-electricdipolesetup, def:physics:phyx-mini-0994:phyxminiproblems-problemphyxmini0994-sourceposition}
  Displacement from a source charge to an arbitrary field point, in metres
\end{definition}
\begin{proof}
  This is the declared model definition of displacement From Source In Meters.
\end{proof}
\begin{definition}[segment Displacement In Meters]
  \label{def:physics:phyx-mini-0994:phyxminiproblems-problemphyxmini0994-segmentdisplacementinmeters}
  \lean{PhyXMiniProblems.ProblemPhyXMini0994.segmentDisplacementInMeters}
  \uses{def:physics:phyx-mini-0994:phyxminiproblems-problemphyxmini0994-figuresegment, def:physics:phyx-mini-0994:phyxminiproblems-problemphyxmini0994-segmentstart, def:physics:phyx-mini-0994:phyxminiproblems-problemphyxmini0994-segmentend, def:physics:phyx-mini-0994:phyxminiproblems-problemphyxmini0994-positionvectorinmeters, def:physics:phyx-mini-0994:phyxminiproblems-problemphyxmini0994-electricdipolesetup}
  Displacement associated with a labelled segment, in metres
\end{definition}
\begin{proof}
  This is the declared model definition of segment Displacement In Meters.
\end{proof}
\begin{definition}[resultant Field At CIn Newtons Per Coulomb]
  \label{def:physics:phyx-mini-0994:phyxminiproblems-problemphyxmini0994-resultantfieldatcinnewtonspercoulomb}
  \lean{PhyXMiniProblems.ProblemPhyXMini0994.resultantFieldAtCInNewtonsPerCoulomb}
  \uses{def:physics:phyx-mini-0994:phyxminiproblems-problemphyxmini0994-electricdipolesetup}
  The total electric-field vector at point c, in newtons per coulomb
\end{definition}
\begin{proof}
  This is the declared model definition of resultant Field At CIn Newtons Per Coulomb.
\end{proof}
\begin{definition}[resultant Field Norm At CIn Newtons Per Coulomb]
  \label{def:physics:phyx-mini-0994:phyxminiproblems-problemphyxmini0994-resultantfieldnormatcinnewtonspercoulomb}
  \lean{PhyXMiniProblems.ProblemPhyXMini0994.resultantFieldNormAtCInNewtonsPerCoulomb}
  \uses{def:physics:phyx-mini-0994:phyxminiproblems-problemphyxmini0994-electricdipolesetup, def:physics:phyx-mini-0994:phyxminiproblems-problemphyxmini0994-resultantfieldatcinnewtonspercoulomb}
  Norm of the total electric-field vector at point c
\end{definition}
\begin{proof}
  This is the declared model definition of resultant Field Norm At CIn Newtons Per Coulomb.
\end{proof}
\begin{definition}[Matches Stated Electric Dipole Scenario]
  \label{def:physics:phyx-mini-0994:phyxminiproblems-problemphyxmini0994-matchesstatedelectricdipolescenario}
  \lean{PhyXMiniProblems.ProblemPhyXMini0994.MatchesStatedElectricDipoleScenario}
  \uses{def:physics:phyx-mini-0994:phyxminiproblems-problemphyxmini0994-lengthinmeters, def:physics:phyx-mini-0994:phyxminiproblems-problemphyxmini0994-chargeinnanocoulombs, def:physics:phyx-mini-0994:phyxminiproblems-problemphyxmini0994-expectedsourcepoint, def:physics:phyx-mini-0994:phyxminiproblems-problemphyxmini0994-expectedchargeinnanocoulombs, def:physics:phyx-mini-0994:phyxminiproblems-problemphyxmini0994-positionvectorinmeters, def:physics:phyx-mini-0994:phyxminiproblems-problemphyxmini0994-electricdipolesetup, def:physics:phyx-mini-0994:phyxminiproblems-problemphyxmini0994-sourceposition}
  The textual problem data: two point charges of equal magnitude and opposite sign, separated by 0.100 m. No electric-field value occurs here
\end{definition}
\begin{proof}
  This is the declared model definition of Matches Stated Electric Dipole Scenario.
\end{proof}
\begin{definition}[Matches Supplied Electric Dipole Figure]
  \label{def:physics:phyx-mini-0994:phyxminiproblems-problemphyxmini0994-matchessuppliedelectricdipolefigure}
  \lean{PhyXMiniProblems.ProblemPhyXMini0994.MatchesSuppliedElectricDipoleFigure}
  \uses{def:physics:phyx-mini-0994:phyxminiproblems-problemphyxmini0994-lengthinmeters, def:physics:phyx-mini-0994:phyxminiproblems-problemphyxmini0994-expectedchargesign, def:physics:phyx-mini-0994:phyxminiproblems-problemphyxmini0994-expectedchargecolor, def:physics:phyx-mini-0994:phyxminiproblems-problemphyxmini0994-expectedsegmentlengthincentimeters, def:physics:phyx-mini-0994:phyxminiproblems-problemphyxmini0994-expectedarroworigin, def:physics:phyx-mini-0994:phyxminiproblems-problemphyxmini0994-expectedarrowdirection, def:physics:phyx-mini-0994:phyxminiproblems-problemphyxmini0994-positionvectorinmeters, def:physics:phyx-mini-0994:phyxminiproblems-problemphyxmini0994-electricdipolesetup, def:physics:phyx-mini-0994:phyxminiproblems-problemphyxmini0994-segmentdisplacementinmeters}
  Primary-image evidence. The predicate records presentation and geometry only; the qualitative arrow transcription is not identified with the independently modeled physical fields
\end{definition}
\begin{proof}
  This is the declared model definition of Matches Supplied Electric Dipole Figure.
\end{proof}
\begin{definition}[Has Physical Electric Dipole Parameters]
  \label{def:physics:phyx-mini-0994:phyxminiproblems-problemphyxmini0994-hasphysicalelectricdipoleparameters}
  \lean{PhyXMiniProblems.ProblemPhyXMini0994.HasPhysicalElectricDipoleParameters}
  \uses{def:physics:phyx-mini-0994:phyxminiproblems-problemphyxmini0994-lengthinmeters, def:physics:phyx-mini-0994:phyxminiproblems-problemphyxmini0994-chargeincoulombs, def:physics:phyx-mini-0994:phyxminiproblems-problemphyxmini0994-electricdipolesetup, def:physics:phyx-mini-0994:phyxminiproblems-problemphyxmini0994-displacementfromsourceinmeters}
  Positivity, nonzero-charge, and source-separation conditions
\end{definition}
\begin{proof}
  This is the declared model definition of Has Physical Electric Dipole Parameters.
\end{proof}
\begin{definition}[Uses School Coulomb Constant]
  \label{def:physics:phyx-mini-0994:phyxminiproblems-problemphyxmini0994-usesschoolcoulombconstant}
  \lean{PhyXMiniProblems.ProblemPhyXMini0994.UsesSchoolCoulombConstant}
  \uses{def:physics:phyx-mini-0994:phyxminiproblems-problemphyxmini0994-electricdipolesetup}
  The rounded Coulomb constant used in the school-level multiple-choice calculation. Physlib's EMSystem.coulombConstant is the scalar 1 / (4 π ε₀) that appears in the Coulomb field law
\end{definition}
\begin{proof}
  This is the declared model definition of Uses School Coulomb Constant.
\end{proof}
\begin{definition}[Satisfies Point Charge Coulomb Field Law]
  \label{def:physics:phyx-mini-0994:phyxminiproblems-problemphyxmini0994-satisfiespointchargecoulombfieldlaw}
  \lean{PhyXMiniProblems.ProblemPhyXMini0994.SatisfiesPointChargeCoulombFieldLaw}
  \uses{def:physics:phyx-mini-0994:phyxminiproblems-problemphyxmini0994-chargeincoulombs, def:physics:phyx-mini-0994:phyxminiproblems-problemphyxmini0994-electricdipolesetup, def:physics:phyx-mini-0994:phyxminiproblems-problemphyxmini0994-sourceposition, def:physics:phyx-mini-0994:phyxminiproblems-problemphyxmini0994-displacementfromsourceinmeters}
  For source charge q and source-to-field-point displacement r, the field is k q r / ‖r‖³. This general all-points law contains no requested result
\end{definition}
\begin{proof}
  This is the declared model definition of Satisfies Point Charge Coulomb Field Law.
\end{proof}
\begin{definition}[Satisfies Electric Field Superposition]
  \label{def:physics:phyx-mini-0994:phyxminiproblems-problemphyxmini0994-satisfieselectricfieldsuperposition}
  \lean{PhyXMiniProblems.ProblemPhyXMini0994.SatisfiesElectricFieldSuperposition}
  \uses{def:physics:phyx-mini-0994:phyxminiproblems-problemphyxmini0994-sourcecharge, def:physics:phyx-mini-0994:phyxminiproblems-problemphyxmini0994-electricdipolesetup}
  The total electric field is the vector sum of the two source fields
\end{definition}
\begin{proof}
  This is the declared model definition of Satisfies Electric Field Superposition.
\end{proof}
\begin{definition}[Resultant Magnitude Represents Field At C]
  \label{def:physics:phyx-mini-0994:phyxminiproblems-problemphyxmini0994-resultantmagnituderepresentsfieldatc}
  \lean{PhyXMiniProblems.ProblemPhyXMini0994.ResultantMagnitudeRepresentsFieldAtC}
  \uses{def:physics:phyx-mini-0994:phyxminiproblems-problemphyxmini0994-fieldmagnitudeinnewtonspercoulomb, def:physics:phyx-mini-0994:phyxminiproblems-problemphyxmini0994-electricdipolesetup, def:physics:phyx-mini-0994:phyxminiproblems-problemphyxmini0994-resultantfieldnormatcinnewtonspercoulomb}
  The dimensionful scalar observable is the Euclidean norm of the independently modeled resultant vector at c; no numerical value is fixed here
\end{definition}
\begin{proof}
  This is the declared model definition of Resultant Magnitude Represents Field At C.
\end{proof}
\begin{definition}[Answer Choice]
  \label{def:physics:phyx-mini-0994:phyxminiproblems-problemphyxmini0994-answerchoice}
  \lean{PhyXMiniProblems.ProblemPhyXMini0994.AnswerChoice}
  Labels attached to the four multiple-choice answers
\end{definition}
\begin{proof}
  This is the declared model definition of Answer Choice.
\end{proof}
\begin{definition}[field Magnitude In Newtons Per Coulomb]
  \label{def:physics:phyx-mini-0994:phyxminiproblems-problemphyxmini0994-answerchoice-fieldmagnitudeinnewtonspercoulomb}
  \lean{PhyXMiniProblems.ProblemPhyXMini0994.AnswerChoice.fieldMagnitudeInNewtonsPerCoulomb}
  \uses{def:physics:phyx-mini-0994:phyxminiproblems-problemphyxmini0994-fieldmagnitudeinnewtonspercoulomb, def:physics:phyx-mini-0994:phyxminiproblems-problemphyxmini0994-answerchoice}
  Electric-field magnitude in newtons per coulomb printed by each choice
\end{definition}
\begin{proof}
  This is the declared model definition of field Magnitude In Newtons Per Coulomb.
\end{proof}
\begin{definition}[recorded Dataset Answer]
  \label{def:physics:phyx-mini-0994:phyxminiproblems-problemphyxmini0994-recordeddatasetanswer}
  \lean{PhyXMiniProblems.ProblemPhyXMini0994.recordedDatasetAnswer}
  \uses{def:physics:phyx-mini-0994:phyxminiproblems-problemphyxmini0994-answerchoice}
  Answer label recorded by the supplied dataset metadata
\end{definition}
\begin{proof}
  This is the declared model definition of recorded Dataset Answer.
\end{proof}
\begin{definition}[Is Rightward]
  \label{def:physics:phyx-mini-0994:phyxminiproblems-problemphyxmini0994-isrightward}
  \lean{PhyXMiniProblems.ProblemPhyXMini0994.IsRightward}
  A vector in the diagram plane points strictly rightward
\end{definition}
\begin{proof}
  This is the declared model definition of Is Rightward.
\end{proof}
\begin{definition}[Rounds To Nearest Hundred]
  \label{def:physics:phyx-mini-0994:phyxminiproblems-problemphyxmini0994-roundstonearesthundred}
  \lean{PhyXMiniProblems.ProblemPhyXMini0994.RoundsToNearestHundred}
  A field magnitude rounds to the displayed nearest 100 N/C
\end{definition}
\begin{proof}
  This is the declared model definition of Rounds To Nearest Hundred.
\end{proof}
\begin{definition}[Is Unique Nearest Displayed Field Magnitude]
  \label{def:physics:phyx-mini-0994:phyxminiproblems-problemphyxmini0994-isuniquenearestdisplayedfieldmagnitude}
  \lean{PhyXMiniProblems.ProblemPhyXMini0994.IsUniqueNearestDisplayedFieldMagnitude}
  \uses{def:physics:phyx-mini-0994:phyxminiproblems-problemphyxmini0994-fieldmagnitudeinnewtonspercoulomb, def:physics:phyx-mini-0994:phyxminiproblems-problemphyxmini0994-answerchoice}
  A displayed choice is strictly nearer than every alternative
\end{definition}
\begin{proof}
  This is the declared model definition of Is Unique Nearest Displayed Field Magnitude.
\end{proof}
\begin{lemma}[resultant Field At C eq coulomb Sum]
  \label{lem:physics:phyx-mini-0994:phyxminiproblems-problemphyxmini0994-resultantfieldatc-eq-coulombsum}
  \lean{PhyXMiniProblems.ProblemPhyXMini0994.resultantFieldAtC_eq_coulombSum}
  \uses{def:physics:phyx-mini-0994:phyxminiproblems-problemphyxmini0994-chargeincoulombs, def:physics:phyx-mini-0994:phyxminiproblems-problemphyxmini0994-sourcecharge, def:physics:phyx-mini-0994:phyxminiproblems-problemphyxmini0994-electricdipolesetup, def:physics:phyx-mini-0994:phyxminiproblems-problemphyxmini0994-displacementfromsourceinmeters, def:physics:phyx-mini-0994:phyxminiproblems-problemphyxmini0994-resultantfieldatcinnewtonspercoulomb, def:physics:phyx-mini-0994:phyxminiproblems-problemphyxmini0994-hasphysicalelectricdipoleparameters, def:physics:phyx-mini-0994:phyxminiproblems-problemphyxmini0994-satisfiespointchargecoulombfieldlaw, def:physics:phyx-mini-0994:phyxminiproblems-problemphyxmini0994-satisfieselectricfieldsuperposition}
  Coulomb's vector law and superposition specialize at c to the sum of the two source expressions. The right-hand side uses only source data and the independent electromagnetic-system constant
\end{lemma}
\begin{proof}
  The conclusion follows from the governing relations and figure readouts named in the statement of resultant Field At C eq coulomb Sum.
\end{proof}
\begin{lemma}[resultant Field At C direction and numerical Bounds]
  \label{lem:physics:phyx-mini-0994:phyxminiproblems-problemphyxmini0994-resultantfieldatc-direction-and-numericalbounds}
  \lean{PhyXMiniProblems.ProblemPhyXMini0994.resultantFieldAtC_direction_and_numericalBounds}
  \uses{def:physics:phyx-mini-0994:phyxminiproblems-problemphyxmini0994-electricdipolesetup, def:physics:phyx-mini-0994:phyxminiproblems-problemphyxmini0994-resultantfieldatcinnewtonspercoulomb, def:physics:phyx-mini-0994:phyxminiproblems-problemphyxmini0994-resultantfieldnormatcinnewtonspercoulomb, def:physics:phyx-mini-0994:phyxminiproblems-problemphyxmini0994-matchesstatedelectricdipolescenario, def:physics:phyx-mini-0994:phyxminiproblems-problemphyxmini0994-matchessuppliedelectricdipolefigure, def:physics:phyx-mini-0994:phyxminiproblems-problemphyxmini0994-hasphysicalelectricdipoleparameters, def:physics:phyx-mini-0994:phyxminiproblems-problemphyxmini0994-usesschoolcoulombconstant, def:physics:phyx-mini-0994:phyxminiproblems-problemphyxmini0994-satisfiespointchargecoulombfieldlaw, def:physics:phyx-mini-0994:phyxminiproblems-problemphyxmini0994-satisfieselectricfieldsuperposition, def:physics:phyx-mini-0994:phyxminiproblems-problemphyxmini0994-isrightward}
  The equal-and-opposite charges and symmetric 5-12-13 geometry make the vertical components cancel and the horizontal components add. With the school Coulomb constant, the resulting norm lies between 4850 N/C and 4950 N/C
\end{lemma}
\begin{proof}
  The conclusion follows from the governing relations and figure readouts named in the statement of resultant Field At C direction and numerical Bounds.
\end{proof}
% --- Archon declaration topology end ---
% --- Archon physics formalization source end ---
